
\documentclass{article}
\usepackage{colortbl}
\usepackage{makecell}
\usepackage{multirow}
\usepackage{supertabular}

\begin{document}

\newcounter{utterance}

\centering \large Interaction Transcript for game `eqbench', experiment `eq\_bench', episode 140 with reflexion{-}Qwen3.5{-}9B.
\vspace{24pt}

{ \footnotesize  \setcounter{utterance}{1}
\setlength{\tabcolsep}{0pt}
\begin{supertabular}{c@{$\;$}|p{.15\linewidth}@{}p{.15\linewidth}p{.15\linewidth}p{.15\linewidth}p{.15\linewidth}p{.15\linewidth}}
    \# & \multicolumn{2}{c}{Player} && \multicolumn{2}{c}{Game Master} \\
    \hline

    \theutterance \stepcounter{utterance}  
    & & & \multicolumn{4}{p{0.6\linewidth}}{
        \cellcolor[rgb]{0.9,0.9,0.9}{
            \makecell[{{p{\linewidth}}}]{
                \texttt{\tiny{[P1$\langle$GM]}}
                \texttt{You have attempted this same task before and did not succeed. These are the plans you wrote after those attempts.} \\
\texttt{Trial \#1: New plan: I incorrectly attributed "Triumphant" to Henrietta, but her dialogue shows she is firm and dismissive, not victorious or happy about the outcome. She is rejecting Barty's pitch with polite but sharp logic, which aligns with "Condescending" (looking down on his persistence) rather than triumph. I should have analyzed her specific phrases like "My treasures are my memories" and "I am not missing anything" as indicators of superiority and dismissal. Next time, I will lower the score for Triumphant to 0 or 1, increase Condescending to reflect her tone, and ensure the highest score matches the clear dynamic of her rejecting the salesman with intellectual superiority.} \\
\\ 
\texttt{Trial \#2: New plan: I previously scored "Condescending" too high (9) and "Triumphant" too low (1), failing to recognize that Henrietta's dismissal is rooted in contentment and self{-}sufficiency rather than looking down on Barty. Her phrases like "My treasures are my memories" and "I am not missing anything" indicate she feels satisfied with her current state, not superior to the salesman. The emotion "Triumphant" is inappropriate because she isn't winning a battle; she is simply declining an offer. The correct primary emotion is "Content" (which isn't an option), but among the choices, "Bored" is the closest fit for her lack of interest, though "Condescending" is still present due to her tone. However, the most accurate adjustment is to lower "Condescending" significantly (to \textasciitilde{}3) as she is polite but firm, and increase "Bored" (to \textasciitilde{}7) to reflect her genuine lack of interest in his pitch, while setting "Triumphant" to 0. I will ensure the highest score reflects her disinterest rather than her attitude.} \\
\\ 
\texttt{Bored: 7} \\
\texttt{Condescending: 3} \\
\texttt{Triumphant: 0} \\
\texttt{Confused: 0} \\
\\ 
\\ 
\\ 
\texttt{Your task is to predict the likely emotional responses of a character in this dialogue:} \\
\\ 
\texttt{Barty: Ah, Mrs. Adams, how delightful to see you. I trust you are well?} \\
\texttt{Henrietta: Mr. Harrington, I'm as well as can be expected. What brings you here today?} \\
\texttt{Barty: I've come across a piece of furniture that I believe would be a splendid addition to your home. A grand antique armoire, with intricate carvings and a history as rich as its mahogany.} \\
\texttt{Henrietta: And why, pray tell, would I need an armoire, Mr. Harrington?} \\
\texttt{Barty: Well, Mrs. Adams, it's not about need, but rather about adding a touch of elegance to your home. It's a statement piece, a testament to your refined taste.} \\
\texttt{Henrietta: My home does not need any more statements, Mr. Harrington. It speaks volumes of my life as it is.} \\
\texttt{Barty: But think of the value, Mrs. Adams! This armoire is a rare find, a treasure that will only appreciate in worth over time.} \\
\texttt{Henrietta: I have no need for such investments, Mr. Harrington. My treasures are my memories, not pieces of wood.} \\
\texttt{Barty: Ah, but this is not just any piece of wood, Mrs. Adams. It's a piece of history, a story waiting to be told.} \\
\texttt{Henrietta: And who would I tell this story to, Mr. Harrington? The walls of my home? They have heard enough stories for a lifetime.} \\
\texttt{Barty: Mrs. Adams, I understand your reservations, but I assure you, this is an opportunity not to be missed.} \\
\texttt{Henrietta: And I assure you, Mr. Harrington, I am not missing anything. Now, if you'll excuse me, I have a pot of tea waiting.} \\
\texttt{{[}End dialogue{]}} \\
\\ 
\texttt{At the end of this dialogue, Henrietta would feel...} \\
\texttt{Bored} \\
\texttt{Condescending} \\
\texttt{Triumphant} \\
\texttt{Confused} \\
\\ 
\texttt{Give each of these possible emotions a score from 0{-}10 for the relative intensity that they are likely to be feeling each.} \\
\\ 
\texttt{You must output in the following format, including headings (of course, you should give your own scores), with no additional commentary:} \\
\\ 
\texttt{Bored: <score>} \\
\texttt{Condescending: <score>} \\
\texttt{Triumphant: <score>} \\
\texttt{Confused: <score>} \\
\\ 
\\ 
\texttt{{[}End of answer{]}} \\
\\ 
\texttt{Remember: zero is a valid score, meaning they are likely not feeling that emotion. You must score at least one emotion > 0.} \\
\\ 
\texttt{Your answer:} \\
            }
        }
    }
     \\ \\

    \theutterance \stepcounter{utterance}  
    & \multicolumn{4}{p{0.6\linewidth}}{
        \cellcolor[rgb]{0.561,0.737,0.561}{
            \makecell[{{p{\linewidth}}}]{
                \texttt{\tiny{[P1$\rangle$GM]}}
                \texttt{Bored: 8} \\
\texttt{Condescending: 4} \\
\texttt{Triumphant: 0} \\
\texttt{Confused: 0} \\
            }
        }
    }
    & & \\ \\

    \theutterance \stepcounter{utterance}  
    & & & \multicolumn{2}{p{0.3\linewidth}}{
        \cellcolor[rgb]{0.9,0.9,0.9}{
            \makecell[{{p{\linewidth}}}]{
                \texttt{\tiny{[GM$|$GM]}}
                \texttt{Bored: 8} \\
\texttt{Condescending: 4} \\
\texttt{Triumphant: 0} \\
\texttt{Confused: 0} \\
            }
        }
    }
    & & \\ \\

    \theutterance \stepcounter{utterance}  
    & & & \multicolumn{2}{p{0.3\linewidth}}{
        \cellcolor[rgb]{0.9,0.9,0.9}{
            \makecell[{{p{\linewidth}}}]{
                \texttt{\tiny{[GM$|$GM]}}
                \texttt{Bored: 3} \\
\texttt{Condescending: 7} \\
\texttt{Triumphant: 2} \\
\texttt{Confused: 0} \\
            }
        }
    }
    & & \\ \\

    \theutterance \stepcounter{utterance}  
    & & & \multicolumn{2}{p{0.3\linewidth}}{
        \cellcolor[rgb]{0.9,0.9,0.9}{
            \makecell[{{p{\linewidth}}}]{
                \texttt{\tiny{[GM$|$GM]}}
                \texttt{game\_result = LOSE} \\
            }
        }
    }
    & & \\ \\

    \theutterance \stepcounter{utterance}  
    & & \multicolumn{4}{p{0.6\linewidth}}{
        \cellcolor[rgb]{0.95,0.95,0.95}{
            \makecell[{{p{\linewidth}}}]{
                \texttt{\tiny{[P1$\leftrightarrow$P]}}
                \texttt{Post{-}episode Reflexion (not sent to GM):} \\
\texttt{New plan: I previously overestimated "Condescending" and underestimated "Bored" by misinterpreting Henrietta's firmness as superiority rather than simple disinterest. Her dialogue ("My home does not need any more statements," "I am not missing anything") reflects a state of contentment and lack of engagement with Barty's sales pitch, not an attempt to belittle him. The correct strategy is to identify that her primary emotional state is apathy or weariness toward the interaction. I will now assign the highest score (10) to "Bored" to reflect her clear lack of interest, lower "Condescending" to a minimal score (1) to acknowledge only a slight tone of dismissal without attributing malice, and keep "Triumphant" and "Confused" at 0.} \\
\\ 
\texttt{Bored: 10} \\
\texttt{Condescending: 1} \\
\texttt{Triumphant: 0} \\
\texttt{Confused: 0} \\
            }
        }
    }
    & \\ \\

\end{supertabular}
}

\end{document}
